\chapter*{Resumen}
\addcontentsline{toc}{chapter}{Resumen}
La pérdida o degradación de canales en señales electroencefalográficas (EEG) constituye un problema recurrente en estudios clínicos, experimentos de neurociencia cognitiva e interfaces cerebro-computador. La ausencia de sensores altera la topografía espacial de la actividad cerebral, dificulta la preservación de potenciales relacionados a eventos y puede introducir sesgos en análisis espectrales posteriores. Esta tesis estudia el problema de reconstrucción de canales EEG faltantes desde el marco del Procesamiento de Señales en Grafos (Graph Signal Processing, GSP), comparando métodos geométricos clásicos, métodos GSP instantáneos y métodos con regularización temporal.

La contribución central es metodológica y experimental: se construye un benchmark reproducible, multi-dataset y multi-métrica, en el cual todos los métodos evaluados, incluyendo los baselines clásicos, son sometidos a optimización sistemática de hiperparámetros. La evaluación considera escenarios de pérdida aleatoria y pérdida espacialmente contigua, niveles de degradación desde pocos canales hasta razones de pérdida severas, y métricas de amplitud, morfología temporal y fidelidad espectral. Los resultados muestran que la regularización temporal, representada principalmente por TRSS, ofrece ventajas consistentes en fidelidad morfológica y espectral bajo el protocolo evaluado, especialmente cuando la información espacial disponible es escasa o está concentradamente degradada. Al mismo tiempo, la tesis enfatiza que la comparación justa reduce la brecha aparente entre algoritmos modernos y baselines clásicos, por lo que las afirmaciones de desempeño deben formularse con trazabilidad cuantitativa y cautela metodológica.

\textbf{Palabras clave:} EEG, interpolación de canales, procesamiento de señales en grafos, GSP, TRSS, Optuna, reconstrucción, regularización temporal, métricas espectrales.
